\documentclass[11pt]{article}

\usepackage[margin=1in]{geometry}
\usepackage{amsmath, amssymb}
\usepackage{booktabs}

\setlength{\parindent}{0pt}
\setlength{\parskip}{4pt}

\newcommand{\code}[1]{\texttt{#1}}

\title{Fixing the Red-Noise Fits:\\
       A Power-Law GP Kernel and a Peak-Masked Continuum}
\author{\texttt{GenerateTable\_CGW}}
\date{2026 September 16}

\begin{document}
\maketitle

\begin{abstract}
The pipeline fits red noise twice: once as a power-law continuum to the Lomb--Scargle
periodogram (the Vaughan fit, \code{rn\_alpha}, \code{rn\_log10\_N}), which sets the
peak-detection threshold that seeds the GP, and once as a kernel term inside the GP.
The synthetic injection--recovery set showed both are broken. The periodogram fit is
biased steep by the periodic signals it is supposed to sit underneath, and the GP term
--- a fixed-$Q$ SHO --- cannot represent a power law at all: its effective slope is
uncorrelated with the true one. This note documents the two fixes: (i) a GP red-noise
term built from a fixed basis of Ornstein--Uhlenbeck (OU) kernels whose weights are tied
to a fitted slope $\alpha$, keeping two free parameters; and (ii) an iteratively
peak-masked continuum fit. On the synthetic set the GP red-noise slope goes from
uncorrelated with the truth ($r=-0.07$, median error 2.09) to $r=0.73$ (median error 0.38),
and the number of detectable-but-missed oscillators falls from 157 to 42.
\end{abstract}

% ===========================================================================
\section{The two problems}
% ===========================================================================

\subsection{The periodogram continuum is biased by the peaks}

\code{\_fit\_red\_noise\_vaughan} regresses $\log_{10}P$ on $\log_{10}f$ over
\emph{every} bin of the periodogram, then adds the Exp(1) log-bias
$-\mathbb{E}[\log_{10}X] = \gamma_E/\ln 10 = 0.25068$. Nothing removes the peaks. Injected
periods are uniform in $P \in [0.1, 9]$\,d, which concentrates them at low frequency,
so the excess power steepens the fitted slope. On the synthetic MW set (246 rows,
\code{02\_compare\_to\_truth.ipynb} \S4.5):

\begin{center}
\begin{tabular}{lcccc}
\toprule
injected oscillators & 0 & 1--3 & 4--6 & 7--10 \\
\midrule
median $\alpha_\mathrm{fit} - \alpha_\mathrm{inj}$ & $-0.146$ & $+0.170$ & $+0.453$ & $+0.712$ \\
\bottomrule
\end{tabular}
\end{center}

This propagates: a steep continuum raises the $\gamma C(f)$ detection threshold at low
frequency, so real low-frequency peaks are never seeded into the GP.

\subsection{The GP red-noise SHO has no slope}

The GP red-noise term was \code{quasisep.SHO} with $Q$ fixed at $1/\sqrt{2}$,
free in $(\sigma, P_\mathrm{rn})$. Its PSD is a broken power law: flat below
$1/P_\mathrm{rn}$ and $\propto f^{-4}$ above. The only freedom is the break position,
so a least-squares slope over $0.1$--$12$\,d$^{-1}$ is a one-to-one function of
$P_\mathrm{rn}$ and saturates at 4. On the synthetic set, that effective slope had
correlation $-0.07$ with the injected $\alpha$ and sat $\approx +2.1$ above it. It is not
a biased estimate of $\alpha$; it carries no information about it.

% ===========================================================================
\section{Fix 1: a power-law GP red-noise kernel}
% ===========================================================================

\subsection{Construction}

A power law is approximated by a sum of Lorentzians with \emph{fixed}, log-spaced break
frequencies $f_j$ and weights tied to $\alpha$:
\begin{equation}
S_\mathrm{RN}(f) = \sum_{j=1}^{N} \frac{c_j}{1 + (f/f_j)^2},
\qquad c_j = K\, f_j^{-\alpha}.
\end{equation}
In the continuum limit, $\int f'^{-\alpha}\,[1+(f/f')^2]^{-1}\,d\ln f' \propto f^{-\alpha}$
(substitute $u = f/f'$; the $u$-integral converges for $0<\alpha<2$), so the sum is a power
law away from the basis edges.

Each Lorentzian is a \code{quasisep.Exp} (OU) kernel,
$k_j(\tau) = \sigma_j^2 e^{-|\tau|/\ell_j}$ with $\ell_j = 1/(2\pi f_j)$. Its one-sided
PSD, in the same convention as \code{sho\_psd} ($\int_0^\infty S\,df = \sigma^2$), is
\begin{equation}
S_j(f) = \frac{2\sigma_j^2}{\pi f_j}\,\frac{1}{1+(f/f_j)^2},
\end{equation}
verified by numerically Fourier-transforming the kernel. Hence
$\sigma_j^2 = \tfrac{\pi}{2} K f_j^{1-\alpha}$. The overall level is set by the PSD at a
pivot $f_0 = 1$\,d$^{-1}$:
\begin{equation}
K = \frac{S_0}{\sum_j f_j^{-\alpha}\,/\,[1 + (f_0/f_j)^2]},
\qquad\text{so that } S_\mathrm{RN}(f_0) = S_0.
\end{equation}
The free parameters are $\theta_\mathrm{RN} = (\ln S_0,\ \alpha)$ --- two, as for the
SHO, whatever $N$ is --- so BIC parameter counts are unchanged. Everything is closed-form
and differentiable, so the JAX gradient path is unaffected.

\textbf{Why $S_0$ and not the total variance.} An earlier prototype
(\code{rn\_kernel\_demo/}) normalised to the total process variance
$\sum_j\sigma_j^2$. For $\alpha<1$ that variance is dominated by the highest basis
terms, which sit above the 24\,d$^{-1}$ Nyquist frequency and are indistinguishable from
white noise, so the amplitude parameter would be degenerate with the jitter. $S_0$ is
set by in-band power only. (That prototype's \code{rn\_psd} was also a factor of 2 low;
slopes were unaffected.)

\subsection{Basis layout}

Accuracy is limited by truncation at the basis edges, not by term density. At fixed
padding, adding terms does not help; extending the basis does, and at 1.2 decades of
padding, 1.5 terms per decade is as good as 6. The default is 1.2 decades beyond each end
of the fitted band at 2 terms per decade: for $0.1$--$12$\,d$^{-1}$ this gives $N=9$,
$f_j \in [6.3\times10^{-3},\,190]$\,d$^{-1}$.

In-band accuracy of that basis (half peak-to-peak log deviation from a pure power law, and
the least-squares slope of the summed PSD over the band):

\begin{center}
\begin{tabular}{lccccc}
\toprule
$\alpha$ (parameter) & 0.25 & 0.5 & 1.0 & 1.5 & 1.75 \\
\midrule
ripple (dex)          & 0.087 & 0.034 & 0.004 & 0.034 & 0.087 \\
in-band slope         & 0.329 & 0.530 & 1.000 & 1.470 & 1.671 \\
\bottomrule
\end{tabular}
\end{center}

Truncation pulls the in-band slope toward 1 at the extremes. The parameter $\alpha$ and
the slope of the spectrum it produces therefore differ by up to 0.08 near the bounds, so
both are reported (next section); \code{rn\_alpha\_band} is the one to compare with a
slope measured from data. $\alpha$ is bounded to $[0.05, 1.95]$, inside the convergence
range. Above 2 the lowest basis term dominates and the in-band shape saturates at
$f^{-2}$, so a slope steeper than 2 is not representable. Rows pinned at the upper
bound should be read as ``$\geq 1.95$''.

\subsection{Code changes}

\begin{description}
\item[\code{gp\_pipeline/gp\_fit.py}]
  New \code{rn\_basis\_frequencies}, \code{\_powerlaw\_rn\_sigmas} (JAX),
  \code{powerlaw\_rn\_psd}, \code{\_unpack\_rn\_powerlaw}, and \code{rn\_psd\_from\_fit},
  which returns the red-noise PSD of any fit dict whichever model it used.
  \code{build\_multi\_sho\_gp}, \code{neg\_log\_likelihood}, \code{make\_initial\_theta},
  \code{make\_bounds}, \code{fit\_multi\_sho\_gp} and \code{iterative\_sho\_gp\_fit} take
  \code{rn\_model} $\in$ \{\code{'powerlaw'}, \code{'sho'}\}, default
  \code{'powerlaw'}, plus \code{rn\_pad\_decades}, \code{rn\_per\_decade}. The fit dict
  now records \code{rn\_model}, \code{rn\_basis\_freqs} and \code{rn\_band}. Initial
  values: $\alpha = 1$, and $S_0$ such that the in-band RN variance is
  $(0.3\,\mathrm{std}(y))^2$, the same fraction the SHO seeding uses.
  \code{summarize\_components} emits \code{rn\_alpha}, \code{rn\_alpha\_band},
  \code{rn\_log10\_S0} and \code{rn\_rms\_band} $= (\int_\mathrm{band} S\,df)^{1/2}$.
\item[\code{gp\_pipeline/gp\_fit\_result.py}]
  \code{kspace\_true}, \code{kspace\_compare\_allfits} and \code{\_\_setstate\_\_} route
  through \code{rn\_psd\_from\_fit} / \code{rn\_model}. Pickles without
  \code{rn\_model} rebuild as SHO, so old results still load.
\item[\code{gp\_pipeline/gp\_table\_io.py}]
  New columns \code{gp\_rn\_alpha}, \code{gp\_rn\_alpha\_band}, \code{gp\_rn\_log10\_S0},
  \code{gp\_rn\_rms\_band}, \code{gp\_rn\_basis\_freqs}; \code{gp\_rn\_sigma/period/Q}
  are NaN under the power-law model. \code{recompute\_gp\_psd\_columns} handles both.
\item[\code{gp\_pipeline/visualize\_gp\_afterfit.py}]
  PSD panels use \code{rn\_psd\_from\_fit}; the component table shows
  $\alpha$ and $S_0$ for a power-law RN term.
\end{description}

\subsection{Checks}

\begin{itemize}
\item Kernel vs.\ analytic PSD: $4\int_0^\infty k(\tau)\cos(2\pi f\tau)\,d\tau$ of the built
  kernel matches \code{powerlaw\_rn\_psd} to $<2\times10^{-6}$ relative at
  $f = 0.1, 1, 5$\,d$^{-1}$; $S_\mathrm{RN}(1) = S_0$ exactly.
\item Regression: \code{rn\_model='sho'} reproduces the stored BIC of the Aug-21 synthetic
  table exactly on the rows tested.
\item Cost: a full 10-component fit takes about the same time with 9 OU terms as with one SHO
  (17 vs.\ 15\,s on a 1-sector row, 19 vs.\ 18\,s on a 2-sector row, JIT included).
\end{itemize}

% ===========================================================================
\section{Fix 2: a peak-masked continuum}
% ===========================================================================

\subsection{Procedure}

\code{\_fit\_red\_noise\_vaughan\_masked} iterates:
\begin{enumerate}
\item fit $\log_{10}P = a + b\log_{10}f$ on the unmasked bins, $C(f) = 10^{a-\beta_k} f^{b}$;
\item flag bins with $P > k\,C(f)$, and mask every bin within $\Delta f = 1/T_\mathrm{eff}$
  of a flagged bin (the periodogram's resolution, i.e.\ the peak's main lobe);
\item repeat until the mask stops changing (cap 20 iterations; stop if fewer than 20\% of
  bins would remain).
\end{enumerate}
The default clip is $k = \ln 100$. Under the null, $X = P/C \sim \mathrm{Exp}(1)$ and the
clip removes $e^{-k} = 1\%$ of pure-noise bins. Because it sits at a fixed multiple of $C$
at every frequency, truncation leaves the \emph{slope} unbiased and shifts only the
intercept. That shift is corrected exactly by replacing $0.25068$ with
\begin{equation}
\beta_k = \mathbb{E}[\log_{10}X \mid X<k]
 = \frac{-\gamma_E - E_1(k) - e^{-k}\ln k}{\ln 10\,(1-e^{-k})},
\end{equation}
which tends to $-0.25068$ as $k\to\infty$ ($\beta_{\ln 100} = -0.2607$). With
$k = \infty$ the masked fit reproduces the old fit to machine precision.

\subsection{Choosing $k$ and the mask width}

Tested on the 246 stored synthetic periodograms (median $\alpha_\mathrm{fit} -
\alpha_\mathrm{inj}$; last row: correlation with $\alpha_\mathrm{inj}$):

\begin{center}
\begin{tabular}{lcccccc}
\toprule
 & unmasked & $k=\gamma_{1\%}$ & $k=\gamma_{1\%}$, $2\Delta f$ & $k=\ln 100$ & $k=\ln 100$, $2\Delta f$ & $k=3$ \\
\midrule
0 osc.     & $-0.146$ & $-0.146$ & $-0.146$ & $-0.136$ & $-0.118$ & $-0.093$ \\
1--3 osc.  & $+0.170$ & $+0.066$ & $+0.061$ & $-0.001$ & $-0.010$ & $-0.003$ \\
4--6 osc.  & $+0.453$ & $+0.300$ & $+0.274$ & $+0.196$ & $+0.189$ & $+0.186$ \\
7--10 osc. & $+0.712$ & $+0.556$ & $+0.444$ & $+0.444$ & $+0.447$ & $+0.353$ \\
\midrule
masked fraction & 0 & 0.015 & 0.023 & 0.093 & 0.155 & 0.337 \\
corr.\ with truth & 0.40 & 0.50 & 0.51 & 0.56 & 0.56 & 0.52 \\
\bottomrule
\end{tabular}
\end{center}

The global 1\% gate ($\gamma \approx 12.7$) masks too little. $k=3$ masks a third of the
spectrum for no gain in correlation. Widening the mask to $2/T_\mathrm{eff}$, or to one or
more single-segment widths ($1/L_\mathrm{seg}$, relevant for multi-sector rows), did not
help consistently. $k = \ln 100$ at $1/T_\mathrm{eff}$ is the default.

Two alternatives were rejected. A threshold-free fit to per-log-bin medians (10 or 20
bins, 50th or 25th percentile) is \emph{worse} than no masking under crowding (+0.96 at
7--10 oscillators): the low-frequency log bins span only a few peak widths, so the peaks
set the median there.

\subsection{How much of the bias is removable}

On the synthetic table, masking removes the bias for sparse spectra but leaves $+0.2$ to
$+0.45$ under crowding. Two controlled tests separate estimator failure from test
design.

\textbf{Idealised periodograms} (independent Exp(1) bins times $C(f)+S(f)$, with sinc$^2$
lines at uniform periods, 1-sector grids): masking takes the bias from $+0.40$ to $+0.02$
for 8 lines at $R = S/C \in [10,100]$, and from $+0.10$ to $+0.02$ for 8 sub-threshold
lines at $R\in[1,3]$.

\textbf{Real sampling.} Using the real timestamps and errors of 60 one-sector rows, red
noise from the generator's \code{red\_noise\_series}, and oscillators drawn from the
generator's own amplitude and period distributions, the shift in fitted $\alpha$
\emph{caused by adding oscillators} (same realisation of red noise and white noise with
and without them) is:

\begin{center}
\begin{tabular}{lccc}
\toprule
oscillators & 0 & 4 & 10 \\
\midrule
unmasked & 0 & $+0.110$ & $+0.231$ \\
masked   & 0 & $+0.025$ & $+0.080$ \\
\bottomrule
\end{tabular}
\end{center}

Masking cuts the oscillator-induced bias about fourfold and adds none on its own. The
larger leftover on the synthetic table is specific to that test (a different red-noise
amplitude distribution, calibrated against the unmasked fit on the longest sector) and was
not chased further.

\subsection{Code changes}

\begin{description}
\item[\code{table\_pipeline/lc\_lsp.py}]
  New \code{\_truncated\_exp\_log10\_mean} and \code{\_fit\_red\_noise\_vaughan\_masked}.
  \code{load\_and\_process\_sectors} takes \code{rn\_mask\_peaks=True},
  \code{rn\_mask\_k=ln 100}, and returns \code{(log10\_N, alpha, mask\_frac)}.
  \code{\_fit\_red\_noise\_vaughan} is unchanged.
\item[\code{table\_pipeline/sector\_expansion.py}]
  \code{expand\_table} passes \code{rn\_mask\_peaks}, \code{rn\_mask\_k} through and emits
  \code{rn\_mask\_frac}. \code{rn\_mask\_peaks=False} reproduces the stored Aug-21 values
  exactly (checked on one cluster).
\item[\code{table\_pipeline/lsp\_stats.py}]
  New \code{refit\_red\_noise\_columns(table)}: the same masked fit from the stored
  \code{LSP\_freq}/\code{LSP\_power}/\code{LC\_t}, without rerunning
  \code{expand\_table}. Old values are kept as \code{rn\_alpha\_unmasked},
  \code{rn\_log10\_N\_unmasked}. Rerun \code{add\_lsp\_metrics} afterwards, and the GP
  fits too if they use \code{threshold\_mode='red\_noise'}.
\end{description}

% ===========================================================================
\section{Results on the synthetic set}
% ===========================================================================

\subsection{Runs}

The 246 synthetic MW rows (124 clusters, 1--3 sectors) were fitted three times with the
\code{add\_gp\_fits} settings of \code{generate\_table.ipynb} (\code{max\_components=10},
\code{fit\_full\_only}, \code{threshold\_mode='red\_noise'}, two-phase fitting, period and
$Q$ priors):

\begin{center}
\begin{tabular}{lll}
\toprule
config & GP red-noise term & continuum used to seed peaks \\
\midrule
\code{sho}            & fixed-$Q$ SHO (the Aug-21 table)     & unmasked \\
\code{powerlaw}       & OU basis, \code{rn\_model='powerlaw'} & unmasked \\
\code{powerlaw+mask}  & OU basis                              & masked (\code{refit\_red\_noise\_columns}) \\
\bottomrule
\end{tabular}
\end{center}

\code{powerlaw} vs.\ \code{sho} isolates the kernel; \code{powerlaw+mask} vs.\
\code{powerlaw} isolates the continuum. Each run took $\sim$11\,min on 30 processes.
Tables: \code{data/cluster\_table\_Sep16\_3sectorcombos\_synthetic\_gpPL\{,\_masked\}.pkl}.
Analysis: \code{synth/compare.py} (\code{slope\_recovery\_table},
\code{gp\_recovery\_table}, \code{bic\_comparison}, \code{plot\_slope\_recovery}),
\S5 of \code{02\_compare\_to\_truth.ipynb}, figure
\code{figs/rn\_fix\_slope\_recovery.png}.

\subsection{Slope recovery}

Median fitted $-$ injected $\alpha$. For the GP, the slope is \code{gp\_rn\_alpha\_band}
(power law) or the break-derived effective slope (SHO).

\begin{center}
\begin{tabular}{llcccccc}
\toprule
config & estimator & 0 osc. & 1--3 & 4--6 & 7--10 & med.\ $|\Delta|$ & $r$ \\
\midrule
\code{sho}           & Vaughan & $-0.146$ & $+0.170$ & $+0.453$ & $+0.712$ & 0.287 & 0.40 \\
\code{sho}           & GP      & $+2.373$ & $+2.066$ & $+1.984$ & $+2.039$ & 2.093 & $-0.07$ \\
\code{powerlaw}      & GP      & $+0.208$ & $+0.410$ & $+0.672$ & $+0.494$ & 0.411 & 0.70 \\
\code{powerlaw+mask} & Vaughan & $-0.136$ & $-0.001$ & $+0.195$ & $+0.444$ & 0.196 & 0.56 \\
\code{powerlaw+mask} & GP      & $+0.208$ & $+0.407$ & $+0.535$ & $+0.459$ & 0.383 & 0.73 \\
\bottomrule
\end{tabular}
\end{center}

The GP red-noise term goes from uncorrelated with the truth ($r=-0.07$) to $r = 0.73$, and
its median error drops from 2.09 to 0.38. Unlike the unmasked Vaughan fit, its bias does
not grow steadily with crowding: the SHO signal components absorb the peaks. The masked
continuum removes the bias entirely for 1--3 oscillators and roughly halves it at 4--10.

\subsection{Where the remaining GP slope bias comes from}

The GP slope is still $+0.2$ to $+0.5$ too steep, including with zero oscillators. Split by
injected slope (\code{powerlaw+mask}):

\begin{center}
\begin{tabular}{lcccc}
\toprule
injected $\alpha$ & $\le 0.5$ & 0.5--1.0 & 1.0--1.5 & $>1.5$ \\
\midrule
GP bias               & $+0.69$ & $+0.51$ & $+0.30$ & $+0.14$ \\
masked Vaughan bias   & $+0.75$ & $+0.04$ & $-0.05$ & $-0.03$ \\
GP at upper $\alpha$ bound & 7\% & 6\% & 13\% & 67\% \\
rows & 83 & 55 & 60 & 48 \\
\bottomrule
\end{tabular}
\end{center}

Three separate effects:

\textbf{(a) Shallow injections are oscillator-dominated.} At $\alpha\le0.5$ \emph{both}
estimators are $\sim+0.7$ steep. The generator draws $(\alpha, \log_{10}N)$ jointly with
correlation $+0.87$, so shallow red noise is also faint. Those spectra are dominated by
the oscillator forest, which is weighted to low frequency (uniform in period). The
spectrum really is steeper than the injected red noise, so this is not a fitting error.

\textbf{(b) The generator's 30-min interpolation grid steepens the time-domain slope.}
For $0.5<\alpha<1.5$ the masked Vaughan fit is unbiased but the GP is not. The generator
builds red noise on a 30-min grid and linearly interpolates it onto the timestamps, a
triangle filter that suppresses power toward Nyquist. The time-domain likelihood sees that
up to 24\,d$^{-1}$; the periodogram fit, capped at 12\,d$^{-1}$, mostly does not.
Controlled test: GP fits with red noise only on the real sampling of 8 one-sector rows,
median $\alpha_\mathrm{band}-\alpha$:

\begin{center}
\begin{tabular}{lccc}
\toprule
generation grid & $\alpha=0.5$ & 1.0 & 1.5 \\
\midrule
30 min (as the generator) & $+0.399$ & $+0.335$ & $+0.239$ \\
2 min (near point sampling) & $-0.084$ & $-0.007$ & $-0.003$ \\
\bottomrule
\end{tabular}
\end{center}

The power-law kernel is unbiased for a point-sampled power law. \textbf{But real data is
not point-sampled.} Each 30-min FFI exposure averages the signal (and \code{resample\_lc}
block-averages 10-min and 200-s data into 30-min bins, so every row carries an effective
30-min boxcar). Repeating the test with a 1-min grid averaged over a 30-min window around
each timestamp (8 rows, amplitude normalised after sampling; interpolation re-run in the
same setup for reference):

\begin{center}
\begin{tabular}{lccc}
\toprule
sampling & $\alpha=0.5$ & 1.0 & 1.5 \\
\midrule
point                        & $-0.149$ & $-0.040$ & $-0.008$ \\
30-min exposure average      & $+0.356$ & $+0.379$ & $+0.277$ \\
30-min grid + interpolation  & $+0.701$ & $+0.621$ & $+0.296$ \\
\bottomrule
\end{tabular}
\end{center}

So exposure averaging \emph{alone} biases the GP slope by $\approx+0.3$. This is a
real-data effect, not only a test artefact. (An earlier draft of this note called it
mild; that was wrong.) The periodogram fit, capped at 12\,d$^{-1}$ where the averaging
filter is $\geq 0.81$, is much less affected. The generator's interpolation adds a
further, test-specific bias on top.

\textbf{(c) The upper bound.} 67\% of rows with injected $\alpha>1.5$ are pinned at
$\alpha = 1.95$, where the in-band slope saturates at $\approx1.80$ (the ceiling in
the figure). The injected slopes reach 2.3, beyond what the construction can represent.
If steep red noise matters, the options are more low-frequency padding (which raises the
saturation slope) or adding an $f^{-4}$ SHO-like term; neither has been tried.

\subsection{Effect on oscillator recovery}

Detectability is judged against the \emph{masked} continuum in every column, so
``detectable'' means the same thing throughout. That is why \code{sho}'s detectable
recovery here (0.743) differs from the 0.957 in notebook \S4, which used its own unmasked
continuum and so counted fewer oscillators as detectable (423 vs.\ 611).

\begin{center}
\begin{tabular}{lccc}
\toprule
 & \code{sho} & \code{powerlaw} & \code{powerlaw+mask} \\
\midrule
recovery, all injected            & 0.425 & 0.460 & 0.592 \\
recovery, detectable (611)        & 0.743 & 0.771 & 0.931 \\
recovery, not detectable          & 0.047 & 0.089 & 0.189 \\
detectable but missed             & 157   & 140   & 42 \\
GP components (total)             & 850   & 850   & 1342 \\
component purity                  & 0.759 & 0.784 & 0.733 \\
mean components per row           & 3.46  & 3.46  & 5.46 \\
\bottomrule
\end{tabular}
\end{center}

The kernel change alone improves recovery modestly (+3\% detectable) and purity
slightly. The masked continuum is the bigger lever. Its lower low-frequency threshold
seeds more real peaks, cutting missed detectable oscillators from 140 to 42, at the cost
of $\sim$2 more components per row and a 5-point drop in purity. The $Q$ filter from
notebook \S4.4 is the natural way to recover that purity.

\textbf{BIC.} Rows share their data across configs, so per-row BIC is comparable.
\code{powerlaw} vs.\ \code{sho}: median $\Delta\mathrm{BIC} = -1.4$, better on 55\% of
rows. That is a small likelihood gain, at the same parameter count, from a term that is
finally the right shape. The \code{powerlaw+mask} BIC median ($+7.5$) mixes in different
component counts. On the 93 rows with equal counts it is $-0.7$, better on 55\%.

% ===========================================================================
\section{Generator change: exposure-averaged sampling}
% ===========================================================================

Following \S4.3, \code{synth/generate.py} now builds $x(t)$ on a 1-min grid
(\code{GRID\_CADENCE\_MIN}) and writes, for each real timestamp, the average of $x$ over
that point's exposure $[t-E/2,\,t+E/2]$. $E$ is the sector's native cadence (median time
step: 29.97, 10 or 3.33\,min), so mixed-cadence clusters are handled per sector
(\code{exposure\_average}, \code{sector\_exposure\_days}). The average uses a running
trapezoid integral, $O(N_\mathrm{grid})$. The red-noise amplitude calibration
(\code{\_calibrate\_red\_scale}) applies the same averaging to the reference sector,
so the calibrated $\log_{10}N$ refers to the data as written. Truth statistics use a
point-sampled 30-min subgrid (\code{TRUTH\_CADENCE\_MIN}), so their definitions are
unchanged.

Checks: a $P=0.1$\,d sinusoid averaged over 30\,min is attenuated by 0.9300 against the
analytic $\mathrm{sinc}(E/P) = 0.9301$ (the injected \code{amplitudes} stay intrinsic);
generated files have identical HDU structure, timestamps and \code{flux\_err} to the real
ones; about 0.5\,s per cluster. The existing \code{light\_curves\_synthetic/} and
\code{ground\_truth.pkl} were \emph{not} regenerated, so the synthetic tables above use
the old sampling. Regenerating means rerunning notebook 01 and then the pipeline.

% ===========================================================================
\section{The $\alpha < 2$ cap}
% ===========================================================================

\subsection{Why the OU basis cannot exceed $f^{-2}$}

\textbf{A hard shape limit.} Each OU term has PSD $\propto [1+(f/f_j)^2]^{-1}$, whose
local log-slope $d\ln S_j/d\ln f = -2(f/f_j)^2/[1+(f/f_j)^2]$ lies in $[-2, 0]$ at every
frequency. The log-slope of a positive sum is a weighted average of the terms' log-slopes
(weights $S_j/\sum S_j$), so it also lies in $[-2,0]$. \emph{No} choice of weights, number
of terms or break frequencies makes a sum of OU kernels steeper than $f^{-2}$ anywhere.

\textbf{An identifiability limit below that.} With $c_j\propto f_j^{-\alpha}$, each term
contributes $c_j f_j^2 f^{-2} \propto f_j^{2-\alpha} f^{-2}$ above its break. For
$\alpha<2$ the higher-frequency terms matter and the sum bends to $f^{-\alpha}$. For
$\alpha\ge 2$ the lowest-frequency term dominates, the in-band shape is simply $f^{-2}$,
and changing $\alpha$ no longer changes the spectrum. The likelihood goes flat and
$\alpha$ drifts. Equivalently, the continuum-limit integral
$\int_0^\infty u^{\alpha-1}/(1+u^2)\,du$ diverges at $\alpha=2$. Hence the bound at 1.95.
With finite padding (1.2\,dec), the in-band slope at $\alpha=1.95$ is only $\approx 1.80$.

\textbf{Raising it.} The limit comes from the OU kernel's $f^{-2}$ tail, not from the
construction. A basis of Mat\'ern-3/2 kernels (\code{quasisep.Matern32}), with PSD
$\propto[1+(f/f_j)^2]^{-2}$ and so log-slopes in $[-4,0]$, admits the same tied-weight
construction for $0<\alpha<4$ (the integral $\int u^{\alpha-1}/(1+u^2)^2\,du$ converges
there). The cost is two state dimensions per term instead of one. Not implemented.

\subsection{Does the real data reach it?}

\textbf{Periodogram slopes.} Refitting the continuum with masking on the two June real
tables (\code{refit\_red\_noise\_columns}; the stored values are unmasked):

\begin{center}
\begin{tabular}{llccccc}
\toprule
table & fit & p50 & p95 & p99 & max & $>1.8$ \,/\, $>1.95$ \\
\midrule
3-sector combos (1250) & unmasked & 1.00 & 1.72 & 1.91 & 2.18 & 3.1\% / 0.9\% \\
                       & masked   & 0.74 & 1.59 & 1.81 & 2.10 & 1.1\% / 0.6\% \\
single sector (2204)   & unmasked & 1.00 & 1.74 & 1.90 & 2.18 & 3.2\% / 0.8\% \\
                       & masked   & 0.74 & 1.62 & 1.86 & 2.10 & 1.8\% / 0.5\% \\
\bottomrule
\end{tabular}
\end{center}

Masking lowers the whole distribution (median $1.00\to0.74$). By this measure only
$\sim$1--2\% of rows are steeper than the OU basis can represent. They are mostly MW
(3.5\% of MW rows vs.\ $\le$1.6\% LMC, 0\% SMC), from 5--9 of 348 clusters.

\textbf{GP slopes.} The GP's own slope is what the cap constrains, and it need not agree.
Power-law GP fits with the notebook settings on 30 single-sector rows with masked
$\alpha>1.5$, and on 30 random rows:

\begin{center}
\begin{tabular}{lccccc}
\toprule
rows & pinned at 1.95 & at 0.05 & median Vaughan & median GP (in-band) & corr. \\
\midrule
masked $\alpha>1.5$ & 12/30 & 3/30 & 1.71 & 1.55 & --- \\
random             & 8/30  & 1/30 & 0.89 & 1.44 & $-0.06$ \\
\bottomrule
\end{tabular}
\end{center}

\textbf{What the pins are.} Each of these 60 rows was refitted with $\alpha$ fixed at 1.0
(everything else identical; the number of SHO components was the same in every pair), and
the log-likelihoods compared, $\Delta\ln L = \ln L_\mathrm{free} - \ln L_{\alpha=1}$. The
red-noise strength is measured as the red-noise PSD at 1\,d$^{-1}$ over the white-noise
PSD $2\sigma^2\Delta t$. The 8 pinned random rows split into:

\begin{center}
\begin{tabular}{lcl}
\toprule
kind & rows & evidence \\
\midrule
no red noise at all       & 4 & RN/WN $\approx 0$, $\Delta\ln L = 0$: amplitude driven to zero, $\alpha$ irrelevant \\
genuinely steep           & 2 & $\Delta\ln L = +14,\ +87$: the data prefer 1.95 over 1 \\
weak preference           & 1 & $\Delta\ln L = +0.9$ \\
optimiser failure         & 1 & $\Delta\ln L = -12.8$: the free fit is \emph{worse} than $\alpha=1$ \\
\bottomrule
\end{tabular}
\end{center}

So the real cap-limited rate in the random sample is $\sim$2--3 of 30, not 8. That is
higher than the 1--2\% from periodogram slopes, as expected: exposure averaging adds
$\approx+0.3$ to GP slopes, and 7--8\% of rows have masked periodogram $\alpha > 1.5$.
In the steep sample, 5 of 12 pins are genuinely steep, 1 has no red noise and 6 are
optimiser failures.

\textbf{Lower bound.} Pins at $\alpha = 0.05$ are harmless. A flat red-noise spectrum is
white noise, degenerate with the jitter term, and every such row had
$|\Delta\ln L| < 7$.

\textbf{Two real problems this exposed.}
\begin{enumerate}
\item \emph{White-noise-dominated rows} do not fit $\alpha \to 0$. The jitter already
  describes the white noise, so the fit drives the red-noise amplitude to zero and
  $\alpha$ floats, often to 1.95. For these rows $\alpha$ should be reported as undefined
  (e.g.\ NaN when RN/WN at 1\,d$^{-1}$ is below a threshold), not read as steep.
\item \emph{The $\alpha$ fit gets stuck.} In 4 of 30 random rows and 13 of 30 steep rows,
  the free fit ends with $\ln L$ more than 1 below the $\alpha=1$ fit, which is impossible
  at a true optimum since $\alpha=1$ is inside the bounds. Seven of these are interior, not
  pinned (worst $\Delta\ln L = -152$). L-BFGS-B from $\alpha_0=1$ is leaving the good basin,
  likely during the two-phase fit while the SHO periods are frozen. Candidate fixes:
  starting $\alpha$ at the masked periodogram slope, or refitting from a second starting
  $\alpha$ and keeping the better result. This also affects the synthetic results in
  \S4.
\end{enumerate}

% ===========================================================================
\section{White-noise flag and stuck $\alpha$ fits}
% ===========================================================================

\subsection{Flag: red noise indistinguishable from white noise}

\code{rn\_white\_noise\_features} (called for every power-law fit) stores
\code{rn\_wn\_ratio} $= S_\mathrm{RN}(1\,\mathrm{d}^{-1}) / [2\,\mathrm{median}(\sigma^2 +
\sigma_\mathrm{jit}^2)\,\Delta t]$ and \code{rn\_near\_white\_noise} $=$ ratio $<$
\code{rn\_wn\_flag\_ratio} (default 0.05). The table columns are \code{gp\_rn\_wn\_ratio} and
\code{gp\_rn\_near\_white\_noise}. Where the flag is set, $\alpha$ is undefined. On the 60
real rows, 16 fits were flagged, and in every one the log-likelihood with $\alpha$ free
differed from $\alpha$ fixed at 1 by at most 2.1: the data do not constrain the slope there.
Pins at the lower bound are not a problem: a flat red-noise spectrum is white noise.

\subsection{Stuck fits: seeding does not help, releasing $\alpha$ last does}

Four strategies on the same 60 real single-sector rows, notebook settings, identical
component counts throughout. ``Fail'' means the final $\ln L$ is more than 1 below the fit
with $\alpha$ fixed at 1. ``Below best'' is counted against the best $\ln L$ any variant
reached for that row.

\begin{center}
\begin{tabular}{llcccc}
\toprule
rows & strategy & fail & worst $\Delta\ln L$ & below best by $>1$ & pinned at 1.95 \\
\midrule
random & free, $\alpha_0=1$ (default)              & 4  & $-16.4$  & 11 & 8 \\
       & free, $\alpha_0$ = masked Vaughan slope    & 5  & $-28.1$  & 10 & 12 \\
       & release last, $\alpha_0=1$                 & 0  & $-0.03$  & 6  & 1 \\
       & release last, $\alpha_0$ = Vaughan         & 4  & $-7.3$   & 4  & 0 \\
\midrule
steep  & free, $\alpha_0=1$ (default)              & 13 & $-152.5$ & 17 & 12 \\
       & free, $\alpha_0$ = Vaughan                 & 11 & $-30.9$  & 11 & 14 \\
       & release last, $\alpha_0=1$                 & 0  & $0.00$   & 11 & 2 \\
       & release last, $\alpha_0$ = Vaughan         & 8  & $-22.9$  & 11 & 4 \\
\bottomrule
\end{tabular}
\end{center}

\textbf{Seeding} with the measured slope (\code{rn\_alpha\_init='measured'}) rescues some
rows (up to +152) but breaks as many (down to $-28$), and 8 failures are unchanged. The
starting $\alpha$ is not the cause. With $\alpha$ held fixed, the oscillator terms and jitter
settle into a better solution; with $\alpha$ free from the start, the joint fit drifts into
a worse one.

\textbf{Releasing $\alpha$ last} (\code{rn\_alpha\_release\_last=True}) holds $\alpha$ at its
seed through the usual stages, then frees it in one extra L-BFGS-B stage started from that
solution. Every stage is kept only if it does not worsen the objective, so the result can
never be worse than the fixed-$\alpha$ fit. With $\alpha_0 = 1$ that guarantee holds (zero
failures), and pins at 1.95 drop from 8 to 1 (random) and 12 to 2 (steep).

\textbf{Correction: that drop is mostly $\alpha$ not moving.} In the final stage, $\alpha$
stays within 0.02 of its seed in 15/30 random and 24/30 steep real rows, against 1/30 and
0/30 for the free fit. The fewer pins came from the release stage stopping in the local
optimum around $\alpha = 1$, not from better fits.

\textbf{Synthetic rerun.} With release-last as the default, the full synthetic comparison
(\code{+release}, same 246 rows, masked continuum) got \emph{worse} than
\code{powerlaw+mask}:

\begin{center}
\begin{tabular}{lcc}
\toprule
 & \code{powerlaw+mask} (free $\alpha$) & \code{+release} \\
\midrule
GP slope: corr.\ with truth / median $|\Delta\alpha|$ & 0.73 / 0.38 & 0.38 / 0.54 \\
GP $\alpha$ pinned at 1.95 / within 0.02 of 1          & 20\% / 2\%  & 3\% / 40\% \\
recovery of detectable oscillators                    & 0.931       & 0.900 \\
detectable but missed                                 & 42          & 61 \\
median $\Delta$BIC vs.\ \code{sho}, equal components  & $-0.7$      & $+5.4$ \\
\bottomrule
\end{tabular}
\end{center}

Per row, $\ln L_\mathrm{release} - \ln L_\mathrm{free}$ has median $-0.01$, but 31\% of rows are
worse by $>1$ and 23\% better by $>1$ (range $-182$ to $+242$). The 99 rows stuck at
$\alpha\approx1$ lose a median 6.3 in $\ln L$. Neither strategy dominates: the free fit
sometimes ends below the fixed-$\alpha$ solution, and release-last usually cannot leave
it. A version that is never worse than either is to run both and keep the higher objective,
at roughly twice the fit cost; not yet tested.

\code{fit\_multi\_sho\_gp}'s two-phase block was rewritten as a stage loop to support this.
With \code{rn\_alpha\_release\_last=False} it reproduces the previous fits exactly
($\Delta\ln L = 0$ on 6 regression rows). The default was switched to \texttt{True} for the
synthetic rerun above and then reverted to \texttt{False} on that result; the run-both
variant was not pursued. Table:
\code{data/cluster\_table\_Sep16\_3sectorcombos\_synthetic\_gpPL\_masked\_release.pkl}.

% ===========================================================================
\section{Regenerated synthetic set: three ways to set the GP $\alpha$}
% ===========================================================================

The synthetic set was regenerated with exposure-averaged sampling (\S5). The seed was
unchanged, so every injected period, amplitude and $(\alpha,\log_{10}N)$ is identical to the
old set; the old truth is kept as \code{ground\_truth\_Aug20\_interp30.pkl}.
\code{expand\_table} (masked continuum) gave the same 367 MW rows, and the notebook's row
selection gave the same 246. Three GP configurations, all power-law RN, masked-continuum
peak seeding, notebook settings:
\code{free} ($\alpha_0=1$, the default); \code{vaughan\_release}
(\code{rn\_alpha\_init='measured'}, \code{rn\_alpha\_release\_last=True}); and
\code{vaughan\_fixed} (\code{rn\_alpha\_fixed=True}, new: $\alpha$ held at the masked
Vaughan slope and not counted in the BIC). The 738 fits took 36\,min on 12 niced workers,
each pinned to one core.
Tables: \code{data/cluster\_table\_Sep17\_3sectorcombos\_synthetic\_expavg\_\{base,gp\_*\}.pkl}.
Figures: \code{sytheticdata\_generate\_compare/figs/rn\_configs\_expavg\_\{1\_scatter,2\_bias,3\_loglike,4\_recovery\}.png}.

Slope statistics are over the 217 rows where the free fit detects red noise
(\code{gp\_rn\_near\_white\_noise} false):

\begin{center}
\begin{tabular}{lccccc}
\toprule
estimator & med.\ bias & med.\ $|\Delta\alpha|$ & $r$ & bias, 0 osc. & bias, 7--10 osc. \\
\midrule
Vaughan (masked)                 & $+0.04$ & 0.20 & 0.55 & $-0.15$ & $+0.35$ \\
GP, $\alpha$ free                & $+0.27$ & 0.27 & 0.71 & $+0.22$ & $+0.31$ \\
GP, Vaughan seed, release last   & $+0.26$ & 0.27 & 0.64 & $+0.21$ & $+0.36$ \\
GP, fixed at Vaughan             & $+0.04$ & 0.19 & 0.55 & $-0.14$ & $+0.30$ \\
\midrule
GP free $-$ constant offset (2-fold CV by cluster) & $+0.01$ & \textbf{0.10} & 0.71 & & \\
GP free, linear calibration (CV)                  & $-0.09$ & 0.16 & 0.68 & & \\
\bottomrule
\end{tabular}
\end{center}

\begin{itemize}
\item \textbf{Regeneration removed the test artefact.} The free GP's bias went from $+0.38$
  to $+0.27$, and to $+0.22$ with no oscillators, which matches the $\approx+0.3$
  expected from 30-min exposure averaging alone.
\item \textbf{Seeding changes nothing once $\alpha$ is released.} It returns to the free
  fit's offset ($+0.26$) with more scatter ($r$ 0.64 vs.\ 0.71). The offset comes from the
  likelihood itself, not the starting point.
\item \textbf{Fixing $\alpha$ at Vaughan} reproduces the Vaughan numbers but fits worse:
  $\ln L$ is more than 1 below the free fit in 73\% of rows (release-last: 26\% worse,
  20\% better, 53\% within $\pm1$). The GP likelihood genuinely prefers a steeper
  $\alpha$, because the kernel lacks the exposure filter.
\item \textbf{Oscillator recovery does not depend on the strategy}: 92.1 / 92.3 /
  91.3\% of detectable oscillators recovered (48 / 47 / 53 missed), purity 73.0 / 72.7 /
  70.9\%.
\item \textbf{The free GP is the best slope estimator after a constant offset}: median
  error 0.10, half the Vaughan fit's, with the best correlation. A linear calibration is
  worse, because the rows pinned at $\alpha=1.95$ (49 of 217; 36 of them with $\geq4$
  oscillators) distort the fit.
\end{itemize}

% ===========================================================================
\section{Correcting the GP slope: constant offset vs.\ exposure-filtered basis}
% ===========================================================================

\subsection{Why the unfiltered basis reads steep}

Every pipeline flux value is an average over its exposure: \code{resample\_lc} puts all rows
on 30-min bins whatever the native cadence, and the native 30-min FFI points are themselves
30-min integrations. A boxcar average of length $E$ multiplies the spectrum by
$\mathrm{sinc}^2(fE)$, which keeps 81\% of the power at 12\,d$^{-1}$ but only 41\% at the
24\,d$^{-1}$ Nyquist frequency. So the data carry
\begin{equation}
S_\mathrm{data}(f) = S_0 (f/f_0)^{-\alpha}\,\mathrm{sinc}^2(fE),
\end{equation}
while the kernel offers only $S_0 (f/f_0)^{-\alpha}$. The likelihood has one way to imitate
the missing roll-off --- tilt the whole power law --- so the fitted $\alpha$ comes out
$\approx0.3$ steep and is an \emph{apparent} slope, part spectrum and part smoothing.

\subsection{Exposure-filtered basis (option 2b), now the default}

\code{iterative\_sho\_gp\_fit(rn\_exposure\_min=30)} --- the default since 2026-09-17 ---
replaces the analytic basis weights with
\code{exposure\_weight\_table}: per $\alpha$ (0.025 grid), log-weight corrections $u_j$ fitted
so that $\sum_j f_j^{-\alpha}e^{u_j}/[1+(f/f_j)^2]$ matches $f^{-\alpha}\,\mathrm{sinc}^2(fE)$
from $5 f_{j,\min}$ to Nyquist (weight 1), and the mean sinc$^2$ envelope up to
$10\times$Nyquist (weight 0.1), with ridge 0.1 and warm starts along $\alpha$. Inside the
likelihood the corrections are linearly interpolated in $\alpha$ (\code{jnp.interp}), so the
kernel is still 9 OU terms:
\begin{equation}
c_j = S_0\,f_0^{\alpha}\,f_j^{-\alpha}\,e^{u_j(\alpha)},
\qquad \sigma_j^2 = \tfrac{\pi}{2} c_j f_j .
\end{equation}
At $\alpha=1$ the corrections are $u \approx (-0.1, 0.0, -0.4, 0.0, -0.6, +0.3, -0.7,
-4.3, -4.6)$: mild adjustments in band, and a $\sim$70-fold suppression of the two terms
above Nyquist, which bends the model's high-frequency end down the way the exposure does.
Cost and parameter count are unchanged --- 712\,s for 246 rows against 722\,s for the
unfiltered fit, still two free parameters, so BIC comparisons still line up. $S_0$ and $\alpha$ now describe
the \emph{intrinsic} spectrum; \code{rn\_alpha\_band} divides the filter back out. The table
matches the filtered power law to 0.02--0.05\,dex over 0.1--12\,d$^{-1}$, and the kernel's
Fourier transform matches its analytic PSD to $<3\times10^{-6}$. New column:
\code{gp\_rn\_exposure\_days}.

Controlled test (8 real one-sector samplings, red noise only, median fitted $-$ true $\alpha$):

\begin{center}
\begin{tabular}{lccc}
\toprule
 & $\alpha=0.5$ & 1.0 & 1.5 \\
\midrule
point samples, unfiltered basis     & $-0.22$ & $-0.04$ & $+0.03$ \\
30-min averages, unfiltered basis   & $+0.35$ & $+0.40$ & $+0.42$ \\
30-min averages, filtered basis     & $+0.03$ & $+0.08$ & $+0.08$ \\
\bottomrule
\end{tabular}
\end{center}

\subsection{Synthetic comparison}

Free $\alpha$ with the filtered basis on the 246 regenerated synthetic rows (246 fits,
12\,min, 12 pinned workers), against the free fit without it. Option 1 subtracts a constant
offset estimated with 2-fold cross-validation \emph{by cluster} (offsets 0.265 and 0.260).
201 rows have red noise detected in both fits. Figures:
\code{figs/rn\_option1\_vs\_2b\_\{1\_scatter,2\_bias,3\_error\_cdf,4\_fit\_and\_recovery\}.png}.

\begin{center}
\begin{tabular}{lcccccc}
\toprule
 & med.\ bias & med.\ $|\Delta\alpha|$ & p84 $|\Delta\alpha|$ & $r$ & bias $\alpha\le0.5$ & bias $\alpha>1.5$ \\
\midrule
Vaughan (masked)              & $+0.01$ & 0.18 & 0.72 & 0.59 & $+0.70$ & $-0.04$ \\
GP free, raw                  & $+0.26$ & 0.26 & 0.47 & 0.78 & $+0.48$ & $+0.17$ \\
option 1: free $-$ constant   & $-0.00$ & 0.08 & 0.33 & 0.78 & $+0.22$ & $-0.10$ \\
option 2b: filtered basis     & $+0.07$ & 0.10 & 0.32 & 0.66 & $+0.19$ & $+0.10$ \\
\bottomrule
\end{tabular}
\end{center}

\begin{itemize}
\item \textbf{Both fix the offset}, with median errors of 0.08 and 0.10 against 0.18 for Vaughan
  and 0.26 raw.
\item \textbf{Option 2b is the better model}: $\ln L$ improves in 64\% of rows (median
  $+2.6$) and worsens in 13\%. Its bias is flatter across $\alpha$, and it has no ceiling:
  option 1 caps at $1.95_{\rm band}-0.26\approx1.53$, which biases steep rows low. It needs
  no constant calibrated on synthetic data. Rows pinned at 1.95 fall from 43 to 20.
\item \textbf{The lower $r$ for 2b comes from 18 outliers} ($|\Delta\alpha|>0.8$): shallow
  injections (median $\alpha$ 0.24) in crowded rows (median 6 oscillators), fitted with the
  maximum 10 GP components, half of them pinned at 1.95. Option 1 compresses the same rows to
  $\approx1.53$, which flatters its $r$. Without them, $r$ is 0.94 (2b) vs.\ 0.92 (option 1),
  and the median error 0.081 vs.\ 0.078. These rows ran out of components, so unmodelled
  oscillator power leaks into the red-noise term; \code{gp\_n\_components == max\_components}
  is a natural quality flag.
\item Oscillator recovery is unchanged (92.1\% vs.\ 92.6\% of detectable; purity 73.0\% vs.\ 73.5\%).
\end{itemize}

\subsection{What is now the default}

\code{rn\_exposure\_min=30.0} in \code{iterative\_sho\_gp\_fit} (and so through
\code{add\_gp\_fits}), with the value also written explicitly in
\code{generate\_table.ipynb}. Option 1 is \emph{not} applied anywhere: no constant is
subtracted from any column. Pass \code{rn\_exposure\_min=None} for point-sampled data, or
the true exposure if the light curves were not resampled --- the filter is only right for
the exposure actually in the data.

Limits to keep in mind: the table is accurate to 0.02--0.05\,dex over 0.1--12\,d$^{-1}$ but
degrades near $\alpha \approx 1.9$ (0.27\,dex at Nyquist), where the OU tails cannot follow a
steep spectrum and the filter cut-off at once; and the 18 crowded, shallow outliers above
are unaffected by any of this, since they come from oscillator power leaking into the
red-noise term once the fit runs out of components.

% ===========================================================================
\section{The white-noise limit: a prior on $\alpha$}
% ===========================================================================

\subsection{The problem}

White noise is the $\alpha \to 0$ limit of red noise, so in principle a vanishing red-noise
term should simply report a flat spectrum. It does not, because the kernel carries a
separate jitter term, and a flat red-noise component is degenerate with it. Faced with a
light curve holding no red noise, the fit therefore drives the red-noise \emph{amplitude}
to zero and leaves $\alpha$ unconstrained.

On the first real-data run (1250 rows) this affected the 417 rows flagged
\code{gp\_rn\_near\_white\_noise}: red noise a factor $10^{-6}$ below the white noise at
1\,d$^{-1}$, an in-band RMS $6\times10^{-4}$ of the flux error --- and a median fitted
$\alpha$ of \textbf{1.76}, with half the rows piled between 1.8 and the 1.95 bound. Read
naively that is a population of steep-spectrum clusters which does not exist, and it is
\emph{anticorrelated} with the red-noise amplitude, so a flexible model could use the
junk $\alpha$ as a proxy for "amplitude is small".

\subsection{The prior}

\code{\_rn\_alpha\_wn\_neg\_log\_prior} adds
\begin{equation}
-\ln p(\alpha) = w_{\rm pr}\,\alpha^2\;\sigma\!\left(-\frac{\log_{10}(r/r_0)}{\Delta}\right),
\qquad r = \frac{S_0}{2\,\mathrm{med}(\sigma_i^2 + \sigma_{\rm jit}^2)\,\Delta t},
\end{equation}
with $\sigma$ the logistic function: $r$ is the red-noise PSD at the pivot over the
white-noise PSD, the same quantity as \code{rn\_wn\_ratio}. Defaults
$w_{\rm pr} = 1$, $r_0 = 0.1$, $\Delta = 0.5$ decades
(\code{rn\_alpha\_wn\_prior\_weight/\_ratio/\_width}; weight 0 disables it). The switch is
sharp enough to leave real red noise alone --- the penalty weight is 1.0 at $r = 10^{-6}$,
0.5 at $r = 0.1$, and 0.04 at $r = 4$, where a typical detected row sits.

Everything it needs is already in $\theta$, so it is differentiable and JIT-compatible
like the period and $Q$ priors. Implementing it meant collapsing the four near-duplicate
prior branches of \code{fit\_multi\_sho\_gp} into a single objective that sums whichever
penalties apply.

\subsection{Synthetic test}

Same 246 rows, with and without the prior, everything else at the defaults.

\begin{center}
\begin{tabular}{lcc}
\toprule
white-noise rows (flagged) & no prior & with prior \\
\midrule
median $\alpha$              & 1.04 & \textbf{0.05} \\
at the 1.95 bound            & 35\% & \textbf{4\%} \\
at the flat limit (0.05)     & 11\% & \textbf{78\%} \\
\midrule
rows with real red noise (154, identical set) & & \\
\midrule
median bias                  & $+0.074$ & $+0.068$ \\
median $|\Delta\alpha|$      & 0.087 & \textbf{0.078} \\
$r$ with truth               & 0.798 & \textbf{0.828} \\
$\alpha$ moved by $>0.05$    & --- & 8\% of rows \\
median $\Delta\ln L$         & --- & $-0.000$ (9\% worse by $>1$) \\
\bottomrule
\end{tabular}
\end{center}

The arbitrary slopes collapse onto the flat limit, and the rows that carry a real slope are
untouched or marginally better. On the white-noise rows $\ln L$ barely moves either
(median 0.000), which is the point: the likelihood is flat in $\alpha$ there, so the prior
picks a value the data do not constrain.

\textbf{State this in any write-up:} $\alpha \approx 0$ in flagged rows is \emph{imposed},
not measured. The defensible reading is "consistent with flat"; \code{gp\_rn\_wn\_ratio} and
\code{gp\_rn\_near\_white\_noise} mark exactly which rows those are.

\subsection{Why not the alternatives}

Dropping the jitter term would also make $\alpha$ continuous through the limit --- white
noise is then just $\alpha = 0$ --- but jitter absorbs underestimated \code{flux\_err}, and
without it that excess is forced into the red-noise term, pulling $\alpha$ low wherever the
errors are optimistic. Substituting the Vaughan slope in flagged rows works (it tends to 0
on a flat spectrum, being an unbounded least-squares fit) but leaves one column with two
definitions and two different biases. Dropping the GP red-noise term altogether is the
worst option: red noise \emph{is} a covariance concept, and without that term the SHO
components absorb the continuum --- the failure mode the $Q$ prior exists to fight.

% ===========================================================================
\section{Summary and follow-ups}
% ===========================================================================

\begin{itemize}
\item \textbf{Defaults changed:} \code{iterative\_sho\_gp\_fit(rn\_model='powerlaw')} and
  \code{expand\_table(rn\_mask\_peaks=True)}. Both old behaviours remain available
  (\code{rn\_model='sho'}, \code{rn\_mask\_peaks=False}) and reproduce old outputs
  exactly. Existing tables: \code{refit\_red\_noise\_columns} $\to$
  \code{add\_lsp\_metrics} $\to$ \code{add\_gp\_fits}.
\item GP red-noise slope: $r=-0.07 \to 0.73$ against the truth, median error $2.09\to0.38$.
  What remains comes from 30-min exposure averaging ($\approx$+0.3, present in real data
  too), the generator's interpolation, and oscillator-dominated shallow injections.
\item Vaughan continuum: crowding bias removed for sparse spectra and about halved for
  crowded ones. The oscillator-induced shift drops about 4$\times$ on real sampling.
\item Oscillator recovery: detectable-but-missed $157 \to 42$.
\item Generator now samples exposure averages from a 1-min grid (\S5); the synthetic set was
  regenerated and rerun (\S8).
\item Real data: masked periodogram slopes exceed 1.8 in 1--2\% of rows. Of the 27\% of
  real GP fits pinned at 1.95, most are rows with no red noise ($\alpha$ undefined) or
  optimiser failures; $\sim$2--3/30 are genuinely cap-limited (\S6.2).
\item In the white-noise limit $\alpha$ is no longer unconstrained: a prior pulls it to the
  flat limit exactly where the red noise is buried (\S10, on by default), which removes the
  spurious pile-up at 1.95 without touching rows that carry a real slope.
\item Exposure-filtered basis weights (\code{rn\_exposure\_min=30}, \S9) are the default:
  the GP slope is unbiased to $+0.07$ with median error 0.10 (raw: $+0.26$, 0.26), and the
  likelihood improves in 64\% of rows.
\item \textbf{Follow-ups:} (1) --- done: the synthetic set was regenerated and rerun;
  (2) --- done: the exposure bias is handled in the kernel rather than by calibration;
  (3) the GP $\alpha$ offset is corrected by the exposure-filtered basis (\S9;
  \code{rn\_exposure\_min=30}, not yet the default); flag rows that hit
  \code{max\_components};
  (4) Mat\'ern-3/2 basis only if the genuinely steep rows matter;
  (5) rerun the real-data table with the new defaults.
\end{itemize}

\end{document}
